\documentclass{article}
\usepackage{amsmath}

\begin{document}

\section*{Operating System Process vs. Firmware}

Let's break down the differences and analogs between a \textbf{running process} in an OS and \textbf{firmware} running on a microcontroller:

\subsection*{1. Operating System Process vs. Firmware}
\begin{itemize}
    \item \textbf{Operating System Process}:
    \begin{itemize}
        \item A \textbf{process} in an operating system is an instance of a running program.
        \item It has its own \textbf{virtual memory} space, including:
        \begin{itemize}
            \item \textbf{Text Segment}: Contains the program’s code (instructions).
            \item \textbf{Data Segment}: Stores global and static variables.
            \item \textbf{Heap}: For dynamically allocated memory during runtime.
            \item \textbf{Stack}: For function calls, local variables, and return addresses.
        \end{itemize}
        \item The OS \textbf{schedules} processes, manages resources, and provides system services.
    \end{itemize}
    \item \textbf{Firmware} (like on your microcontroller):
    \begin{itemize}
        \item Firmware is essentially a \textbf{single, continuous program} that runs directly on the microcontroller.
        \item It does \textbf{not have virtual memory} or the concept of separate processes.
        \item The code runs from \textbf{flash memory}, and the data is either in \textbf{flash}, \textbf{SRAM}, or \textbf{peripheral memory}.
    \end{itemize}
\end{itemize}

\subsection*{2. Flash vs. RAM for Firmware}
\begin{itemize}
    \item \textbf{Flash Memory} (Non-Volatile):
    \begin{itemize}
        \item Holds the \textbf{program code} (instructions) and \textbf{read-only constants}.
        \item When the microcontroller boots, it starts executing code directly from \textbf{flash}.
        \item Flash memory is not typically used for runtime data storage because writing to it is slow and limited.
    \end{itemize}
    \item \textbf{SRAM (Static RAM)}:
    \begin{itemize}
        \item SRAM is the microcontroller's \textbf{volatile memory} used for \textbf{dynamic data}.
        \item Global variables (if mutable) are initialized by copying from flash to SRAM at startup.
        \item The \textbf{stack} and \textbf{heap} for runtime memory management are located in SRAM.
    \end{itemize}
\end{itemize}

\subsection*{3. Interrupt Table (Vector Table) in Flash}
\begin{itemize}
    \item The \textbf{Interrupt Vector Table} is stored in \textbf{flash memory} by default and contains addresses of the interrupt service routines (ISRs).
    \item This table persists in flash since it doesn't change dynamically during runtime.
\end{itemize}

\subsection*{4. Analogies Between an OS Process and Firmware}
\begin{itemize}
    \item \textbf{Text Segment (Process)} $\leftrightarrow$ \textbf{Flash Memory (Firmware)}:
    \begin{itemize}
        \item The code (instructions) in both cases is stored in non-volatile memory (disk for processes, flash for firmware).
    \end{itemize}
    \item \textbf{Data Segment (Process)} $\leftrightarrow$ \textbf{Flash/SRAM (Firmware)}:
    \begin{itemize}
        \item In a process, static and global variables reside in the data segment. In firmware, read-only data stays in flash, and read/write data gets copied to SRAM during initialization.
    \end{itemize}
    \item \textbf{Heap (Process)} $\leftrightarrow$ \textbf{Heap in SRAM (Firmware)}:
    \begin{itemize}
        \item Both use a heap for dynamically allocated memory, but in firmware, it's allocated directly from SRAM.
    \end{itemize}
    \item \textbf{Stack (Process)} $\leftrightarrow$ \textbf{Stack in SRAM (Firmware)}:
    \begin{itemize}
        \item Both use a stack for managing function calls and local variables. In firmware, the stack is located in SRAM.
    \end{itemize}
\end{itemize}

\subsection*{5. Absence of Virtual Memory}
\begin{itemize}
    \item Firmware lacks \textbf{virtual memory} and runs directly in \textbf{physical memory} (flash and SRAM).
    \item There’s no paging or segmentation as in an OS, meaning firmware runs with full access to the hardware's memory map.
\end{itemize}

\subsection*{6. Single Continuous Program vs. Multitasking Processes}
\begin{itemize}
    \item Firmware is typically a \textbf{single continuous program} running on the microcontroller. 
    \item There’s no concept of multiple processes or multitasking unless you implement a \textbf{real-time operating system (RTOS)}.
\end{itemize}

\end{document}
